\section{The $Q_i$ certificate and the Vandiver bridge}\label{sec:certificate}

\begin{definition}[$Q_i$ / Vandiver certificate]\label{def:vcert}
  \lean{FltVandiver.QiCert.vandiverCert}
  \leanok
  The Boolean certificate \code{vandiverCert}~$p\,\ell\,t$ asserting that the
  $Q_i$ cyclotomic-unit test fires at every even index for the auxiliary pair
  $(\ell, t)$.
\end{definition}

\begin{definition}[Fast sliced certificate]\label{def:vcertfast}
  \lean{FltVandiver.QiCert.vandiverCertFast}
  \leanok
  \uses{def:vcert}
  A pure-\code{Nat} reformulation \code{vandiverCertFast} of the certificate, with
  exponents reduced modulo $\ell - 1$, that \code{native\_decide} evaluates for every
  prime below $1000$.
\end{definition}

\begin{definition}[Machine-speed certificate]\label{def:fast2cert}
  \lean{FltVandiver.QiCert.Fast2.fast2Cert}
  \leanok
  \uses{def:vcert}
  The evaluator \code{fast2Cert} of the all-even test by discrete-logarithm sums:
  $Q_i^{\,k} = h^{\,S_i - N d_i}$ with $h = t^k$ of order $p$, so an index costs
  $p/2$ multiplications modulo $p$ once $O(p)$ tables are built. Every table fact it
  uses is re-verified inside the Boolean. Precompiled; it carries the two
  Wolstenholme primes ($16843$ in seconds, $2124679$ in $31$ core-hours).
\end{definition}

\begin{theorem}[Machine-speed certificate $\Rightarrow$ certificate]\label{thm:fast2bridge}
  \lean{FltVandiver.QiCert.Fast2.vandiverCert_of_fast2}
  \leanok
  \uses{def:fast2cert, def:vcert}
  If \code{fast2Cert}~$p\,\ell\,t\,\mathit{idx}$ holds then so does
  \code{vandiverCert}~$p\,\ell\,t\,\mathit{idx}$: the two evaluators agree on what they
  certify, on the three standard axioms.
\end{theorem}

\begin{theorem}[Base independence of the $Q_i$ test]\label{thm:baseindep}
  \lean{FltVandiver.QiCert.qi_pow_eq_one_iff_of_base}
  \leanok
  \uses{def:vcert}
  For two valid bases $t, t'$ and an even index $i < p$, $Q_i(t)^k = 1$ if and only if
  $Q_i(t')^k = 1$: the verdict of the test depends on $(p, \ell)$ alone, so a passing
  certificate at every even index passes for every valid base
  (\code{vandiverCert\_evenIndices\_of\_base}). A witness is the prime $\ell$.
\end{theorem}

\begin{definition}[Unit index {$[E:C]=\hplus$} (external)]\label{ext:unitindex}
  \leanok
  Washington's Theorem~8.2, \code{cyclotomic\_unit\_index} in
  \code{flt-cyclotomic-nt}: the index of the cyclotomic units equals $\hplus$. It
  turns the $Q_i$ test into a statement about $p \nmid \hplus$.
\end{definition}

\begin{theorem}[$Q_i$ certificate $\Rightarrow$ Vandiver]\label{thm:vbridge}
  \lean{FltVandiver.QiCert.qiVandiverBridge_all}
  \leanok
  \uses{def:vcert, ext:unitindex}
  If \code{vandiverCert}~$p\,\ell\,t$ holds then $p \nmid \hplus$: the $Q_i$ test
  firing at every even index rules out each eigen cyclotomic unit being a $p$-th
  power.
\end{theorem}

\begin{definition}[Small-witness certificate]\label{def:swcert}
  \lean{FltVandiver.SmallWitnessCert}
  \leanok
  \uses{def:vcert}
  \code{SmallWitnessCert}~$p$ packages the pair $(\ell, t)$ with $\ell < p^2-p$ and
  the $Q_i$ certificate --- the single finite input that powers both the Vandiver
  conclusion and the Case~II descent.
\end{definition}
